\documentclass[12pt]{article}
\usepackage{amsmath,amssymb,amsthm}
\usepackage[letterpaper,margin=1in]{geometry}
\usepackage{enumitem}
\usepackage{fancyhdr}

\pagestyle{fancy}
\fancyhf{}
\lhead{MATH 221 --- Problem Set 5}
\rhead{Jordan Kim}
\cfoot{\thepage}

\newtheorem*{claim}{Claim}
\newcommand{\R}{\mathbb{R}}

\begin{document}

\begin{center}
  {\Large\bfseries Problem Set 5: Eigenvalues and Diagonalization}\\[4pt]
  Due Friday, October 9
\end{center}

\section*{Problem 1}
Find the eigenvalues and eigenvectors of
\[
A = \begin{bmatrix} 4 & 1 \\ 2 & 3 \end{bmatrix}.
\]

\begin{proof}[Solution]
The characteristic polynomial is
\[
\det(A - \lambda I) = \begin{vmatrix} 4-\lambda & 1 \\ 2 & 3-\lambda \end{vmatrix}
 = \lambda^2 - 7\lambda + 10 = (\lambda-2)(\lambda-5).
\]
For $\lambda_1 = 5$ we solve $(A-5I)\mathbf{v}=\mathbf{0}$ and obtain
$\mathbf{v}_1 = (1,1)^T$. For $\lambda_2 = 2$ we obtain $\mathbf{v}_2 = (1,-2)^T$.
\end{proof}

\section*{Problem 2}
Let $T\colon \R^3 \to \R^3$ be linear with $T^2 = T$.
\begin{enumerate}[label=(\alph*)]
  \item Show that every eigenvalue of $T$ is $0$ or $1$.
  \item Show that $\R^3 = \ker T \oplus \operatorname{im} T$.
\end{enumerate}

\begin{claim}
If $T\mathbf{v} = \lambda\mathbf{v}$ with $\mathbf{v}\neq\mathbf{0}$, then $\lambda\in\{0,1\}$.
\end{claim}
\begin{proof}
We have $\lambda\mathbf{v} = T\mathbf{v} = T^2\mathbf{v} = \lambda^2\mathbf{v}$, so
$\lambda^2 = \lambda$ and therefore $\lambda(\lambda-1)=0$.
\end{proof}

For (b), write $\mathbf{x} = (\mathbf{x} - T\mathbf{x}) + T\mathbf{x}$. The first
summand lies in $\ker T$ because $T(\mathbf{x}-T\mathbf{x}) = T\mathbf{x}-T^2\mathbf{x}=\mathbf{0}$,
and the second lies in $\operatorname{im} T$. If $\mathbf{y}\in\ker T\cap\operatorname{im}T$,
then $\mathbf{y}=T\mathbf{z}$ and $\mathbf{0}=T\mathbf{y}=T^2\mathbf{z}=T\mathbf{z}=\mathbf{y}$.

\section*{Problem 3}
Diagonalize $A$ from Problem 1, i.e.\ find $P$ and $D$ with $A = PDP^{-1}$.
\[
P = \begin{pmatrix} 1 & 1 \\ 1 & -2 \end{pmatrix}, \qquad
D = \begin{pmatrix} 5 & 0 \\ 0 & 2 \end{pmatrix}, \qquad
P^{-1} = \frac{1}{3}\begin{pmatrix} 2 & 1 \\ 1 & -1 \end{pmatrix}.
\]
Hence $A^{n} = P D^{n} P^{-1}$ for every $n \ge 1$.

\end{document}
